\documentclass[12pt,a4paper]{article}
\usepackage[UTF8]{ctex}
\usepackage{amsmath,amssymb,amsthm}
\usepackage{geometry}

\geometry{margin=2.5cm}

\title{公式11.16与11.18的区别}
\author{第11章补充说明}
\date{}

\begin{document}

\maketitle

\section{公式回顾}

\subsection{公式11.16：使用伴随表示}

\begin{equation}
V_b(t) = [Ad_{X^{-1}X_d}] V_d(t) + K_p X_e(t) + K_i \int_0^t X_e(t) dt
\label{eq:11.16}
\end{equation}

其中：
\begin{itemize}
\item $V_b(t)$是体坐标系下的twist（6维速度向量）
\item $V_d(t)$是期望的twist
\item $[Ad_{X^{-1}X_d}]$是伴随表示矩阵
\item $X_e(t)$是位姿误差（6维向量）
\item $K_p$和$K_i$是比例和积分增益矩阵
\end{itemize}

\subsection{公式11.18：使用显式旋转矩阵表示}

\begin{equation}
\begin{bmatrix}
\omega_b(t) \\
\dot{p}(t)
\end{bmatrix}
=
\begin{bmatrix}
R^T(t)R_d(t) & 0 \\
0 & I
\end{bmatrix}
\begin{bmatrix}
\omega_d(t) \\
\dot{p}_d(t)
\end{bmatrix}
+ K_p X_e(t) + K_i \int_0^t X_e(t) dt
\label{eq:11.18}
\end{equation}

其中：
\begin{itemize}
\item $\omega_b(t)$是体坐标系下的角速度（3维向量）
\item $\dot{p}(t)$是体坐标系下的线速度（3维向量）
\item $\omega_d(t)$和$\dot{p}_d(t)$是期望的角速度和线速度
\item $R(t)$是当前旋转矩阵
\item $R_d(t)$是期望旋转矩阵
\item $I$是3×3单位矩阵
\end{itemize}

\section{主要区别}

\subsection{数学表示形式}

\textbf{公式11.16（伴随表示）：}
\begin{itemize}
\item 使用\textbf{伴随表示}（Adjoint representation）$[Ad_{X^{-1}X_d}]$
\item 紧凑的矩阵形式：$6 \times 6$矩阵作用于6维twist向量
\item 数学上更优雅，符合李群理论
\end{itemize}

\textbf{公式11.18（显式表示）：}
\begin{itemize}
\item 使用\textbf{显式的旋转矩阵表示}
\item 将twist分解为角速度$\omega$和线速度$\dot{p}$
\item 更直观，便于理解和实现
\end{itemize}

\subsection{前馈项的表示}

\textbf{公式11.16的前馈项：}
\begin{equation}
[Ad_{X^{-1}X_d}] V_d(t)
\end{equation}

伴随表示矩阵$[Ad_{X^{-1}X_d}]$的完整形式为：
\begin{equation}
[Ad_{X^{-1}X_d}] = \begin{bmatrix}
R^T R_d & 0 \\
[p]^T R^T R_d & R^T R_d
\end{bmatrix}
\end{equation}

其中$[p]$是平移向量$p$的反对称矩阵。

\textbf{公式11.18的前馈项：}
\begin{equation}
\begin{bmatrix}
R^T(t)R_d(t) & 0 \\
0 & I
\end{bmatrix}
\begin{bmatrix}
\omega_d(t) \\
\dot{p}_d(t)
\end{bmatrix}
\end{equation}

这是伴随表示矩阵的\textbf{简化形式}，假设平移项的影响可以忽略或已单独处理。

\subsection{关键差异：平移项的耦合}

\textbf{公式11.16（完整伴随表示）：}

伴随表示矩阵考虑了旋转和平移的耦合：
\begin{equation}
[Ad_{X^{-1}X_d}] = \begin{bmatrix}
R^T R_d & 0 \\
[p]^T R^T R_d & R^T R_d
\end{bmatrix}
\end{equation}

注意下左块$[p]^T R^T R_d$：这表示旋转对平移的影响（科里奥利效应）。

\textbf{公式11.18（简化形式）：}

假设旋转和平移解耦：
\begin{equation}
\begin{bmatrix}
R^T R_d & 0 \\
0 & I
\end{bmatrix}
\end{equation}

下左块为零，意味着忽略了旋转对平移的耦合影响。

\section{详细数学分析}

\subsection{伴随表示的定义}

对于SE(3)群元素$X = \begin{bmatrix} R & p \\ 0 & 1 \end{bmatrix}$，其伴随表示$[Ad_X]$定义为：
\begin{equation}
[Ad_X] = \begin{bmatrix}
R & 0 \\
[p]R & R
\end{bmatrix}
\end{equation}

其中$[p]$是平移向量$p$的反对称矩阵：
\begin{equation}
[p] = \begin{bmatrix}
0 & -p_z & p_y \\
p_z & 0 & -p_x \\
-p_y & p_x & 0
\end{bmatrix}
\end{equation}

\subsection{公式11.16的完整展开}

公式11.16可以展开为：
\begin{align}
V_b(t) &= [Ad_{X^{-1}X_d}] V_d(t) + K_p X_e(t) + K_i \int_0^t X_e(t) dt \\
\begin{bmatrix}
\omega_b \\
v_b
\end{bmatrix} &= \begin{bmatrix}
R^T R_d & 0 \\
[p]^T R^T R_d & R^T R_d
\end{bmatrix}
\begin{bmatrix}
\omega_d \\
v_d
\end{bmatrix}
+ K_p X_e + K_i \int_0^t X_e dt
\end{align}

其中：
\begin{align}
\omega_b &= R^T R_d \omega_d + (K_p X_e + K_i \int X_e dt)_{\omega} \\
v_b &= [p]^T R^T R_d \omega_d + R^T R_d v_d + (K_p X_e + K_i \int X_e dt)_{v}
\end{align}

注意$v_b$中的项$[p]^T R^T R_d \omega_d$：这表示旋转运动对线速度的影响。

\subsection{公式11.18的展开}

公式11.18可以展开为：
\begin{align}
\omega_b &= R^T R_d \omega_d + (K_p X_e + K_i \int X_e dt)_{\omega} \\
\dot{p} &= I \cdot \dot{p}_d + (K_p X_e + K_i \int X_e dt)_{p}
\end{align}

注意：$\dot{p} = v$（线速度），且忽略了旋转对平移的耦合项。

\section{物理意义对比}

\subsection{旋转对平移的耦合}

\textbf{公式11.16（考虑耦合）：}

当末端执行器在旋转时，其平移速度会受到旋转的影响。这是因为：
\begin{itemize}
\item 旋转中心可能不在原点
\item 旋转会产生科里奥利效应
\item 伴随表示矩阵的下左块$[p]^T R^T R_d$捕获了这个效应
\end{itemize}

\textbf{公式11.18（忽略耦合）：}

假设旋转和平移独立，或者：
\begin{itemize}
\item 旋转中心在原点附近
\item 平移速度相对于旋转速度很小
\item 可以忽略耦合项的影响
\end{itemize}

\subsection{适用场景}

\textbf{公式11.16适用于：}
\begin{itemize}
\item 需要精确考虑旋转-平移耦合的情况
\item 末端执行器远离旋转中心
\item 高速旋转运动
\item 需要严格数学正确性的应用
\end{itemize}

\textbf{公式11.18适用于：}
\begin{itemize}
\item 旋转-平移耦合可以忽略的情况
\item 末端执行器接近旋转中心
\item 低速或小角度旋转
\item 需要简化计算的应用
\end{itemize}

\section{计算复杂度对比}

\subsection{公式11.16的计算}

需要计算：
\begin{enumerate}
\item $X^{-1} X_d$：相对变换（4×4矩阵乘法）
\item $[Ad_{X^{-1}X_d}]$：伴随表示矩阵（需要计算$[p]$和矩阵乘法）
\item $[Ad_{X^{-1}X_d}] V_d$：6×6矩阵与6维向量相乘
\end{enumerate}

\textbf{计算量：}约$O(1)$（常数时间，但常数较大）

\subsection{公式11.18的计算}

需要计算：
\begin{enumerate}
\item $R^T R_d$：3×3矩阵乘法
\item 分块矩阵与向量的乘法（更简单）
\end{enumerate}

\textbf{计算量：}约$O(1)$（常数时间，常数较小）

\textbf{结论：}公式11.18计算更简单，因为：
\begin{itemize}
\item 不需要计算反对称矩阵$[p]$
\item 不需要计算完整的伴随表示矩阵
\item 分块矩阵乘法更直接
\end{itemize}

\section{精度对比}

\subsection{理论精度}

\textbf{公式11.16：}
\begin{itemize}
\item 数学上完全正确
\item 考虑了所有SE(3)群的结构
\item 适用于任意旋转和平移
\end{itemize}

\textbf{公式11.18：}
\begin{itemize}
\item 是公式11.16的近似
\item 忽略了旋转-平移耦合项
\item 在小角度或特定条件下才准确
\end{itemize}

\subsection{误差分析}

公式11.18相对于公式11.16的误差为：
\begin{equation}
\text{误差} = [p]^T R^T R_d \omega_d
\end{equation}

这个误差的大小取决于：
\begin{itemize}
\item 平移向量$p$的大小
\item 旋转角度的大小
\item 角速度$\omega_d$的大小
\end{itemize}

\textbf{当误差可忽略时：}
\begin{itemize}
\item $|p|$很小（末端执行器接近原点）
\item 旋转角度很小
\item $\omega_d$很小（低速旋转）
\end{itemize}

\section{选择建议}

\subsection{何时使用公式11.16？}

\begin{itemize}
\item 需要高精度控制
\item 末端执行器远离旋转中心
\item 高速或大角度旋转
\item 需要严格的数学正确性
\item 计算资源充足
\end{itemize}

\subsection{何时使用公式11.18？}

\begin{itemize}
\item 旋转-平移耦合可以忽略
\item 需要简化计算
\item 末端执行器接近旋转中心
\item 低速或小角度运动
\item 计算资源有限
\end{itemize}

\section{总结}

\textbf{核心区别：}

\begin{enumerate}
\item \textbf{数学表示：}
   \begin{itemize}
   \item 公式11.16：使用伴随表示（Adjoint representation）
   \item 公式11.18：使用显式旋转矩阵表示
   \end{itemize}

\item \textbf{旋转-平移耦合：}
   \begin{itemize}
   \item 公式11.16：考虑完整的耦合项$[p]^T R^T R_d \omega_d$
   \item 公式11.18：忽略耦合项，假设解耦
   \end{itemize}

\item \textbf{计算复杂度：}
   \begin{itemize}
   \item 公式11.16：需要计算完整的伴随表示矩阵
   \item 公式11.18：计算更简单，分块矩阵乘法
   \end{itemize}

\item \textbf{精度：}
   \begin{itemize}
   \item 公式11.16：数学上完全正确
   \item 公式11.18：近似形式，在特定条件下准确
   \end{itemize}

\item \textbf{适用场景：}
   \begin{itemize}
   \item 公式11.16：高精度、高速、大角度应用
   \item 公式11.18：简化计算、低速、小角度应用
   \end{itemize}
\end{enumerate}

\textbf{关键点：}公式11.18是公式11.16的简化形式，通过忽略旋转-平移耦合项来简化计算。选择哪个取决于精度要求、计算资源和运动特性。

\textbf{关系：}公式11.18可以看作是公式11.16在特定条件下的近似，当旋转-平移耦合可以忽略时，两者结果相近。

\end{document}
